\chapter{Problem Statement}
\pagenumbering{arabic} %DO NOT REMOVE FROM THE INTRODUCTION CHAPTER
This template has been adapted for TU Delft assignments and reports. Replace
this chapter, and the ones that follow, with your own work --- the structure
below (problem statement, analysis, results, appendix) is a starting point,
not a requirement.

\section{Assignment}
An aluminium 6061-T6 L-bracket cantilevers a \SI{2.4}{\kilo\gram} vibration
sensor a distance $L$ from a vertical mounting plate, as sketched in
\cref{fig:bracket}. The bracket has a rectangular cross-section of width $b$
and thickness $t$, and a bolt hole of diameter $d$ is drilled through the
root section to fix it to the plate. During operation the sensor mount
experiences a cyclic vertical load oscillating between $F_{\text{min}}$ and
$F_{\text{max}}$ at the rig's operating frequency.

\begin{questions}
    \item What is the maximum bending stress at the root section, and does
    the bracket meet the static factor-of-safety requirement of
    \cref{tab:requirements}?
    \begin{questions}
        \item What is the stress-concentration factor $K_t$ at the bolt
        hole, and what is the resulting local peak stress?
    \end{questions}
    \item Using a Goodman diagram, estimate the fatigue life at the bolt
    hole. Does the bracket meet the minimum fatigue-life requirement of
    \cref{tab:requirements}?
    \item Propose one design change that improves the fatigue margin without
    increasing bracket mass by more than 10\%.
\end{questions}

\begin{table}[H]
    \centering
    \caption{Design requirements}
    \label{tab:requirements}
    \begin{tabularx}{\textwidth}{L C C}
        \toprule
        Requirement & Symbol & Value \\
        \midrule
        Static factor of safety & $\mathrm{FoS}_{\text{static}}$ & $\geq 1.5$ \\
        Minimum fatigue life & $N_f$ & $\geq \num{1e6}$ cycles \\
        \hdashline
        Maximum bracket mass & $m$ & $\leq \SI{0.4}{\kilo\gram}$ \\
        \bottomrule
    \end{tabularx}
\end{table}

\section{Free-Body Analysis}
Treating the bracket as a cantilever beam fixed at the mounting plate, the
root bending moment under the peak load $F_{\text{max}}$ is
\begin{equation}
    \label{eq:moment}
    M = F_{\text{max}} \, L \, ,
\end{equation}
and the corresponding nominal bending stress at the outer fibre of the root
section is
\begin{equation}
    \label{eq:stress}
    \sigma_{\text{nom}} = \frac{M \, c}{I} = \frac{6\,F_{\text{max}}\,L}{b\,t^{2}} \, ,
\end{equation}
where $c = t/2$ and $I = b\,t^{3}/12$ for the rectangular section. The
governing symbols are summarised in \cref{tab:exampletable}; the notation
follows the standard beam-theory convention, e.g.\ \citet{crandall2012energy}.

\begin{table}[H]
    \centering
    \caption{Symbols used in the free-body analysis}
    \label{tab:exampletable}
    \begin{tabularx}{\textwidth}{L C L}
        \toprule
        Symbol & Units & Description \\
        \midrule
        $L$ & \unit{\meter} & Cantilever length, plate to sensor mount \\
        $b$ & \unit{\meter} & Bracket cross-section width \\
        $t$ & \unit{\meter} & Bracket cross-section thickness \\
        $d$ & \unit{\meter} & Bolt-hole diameter at root section \\
        $K_t$ & -- & Stress-concentration factor at the bolt hole \\
        \bottomrule
    \end{tabularx}
\end{table}

\section{Figures}
\Cref{fig:bracket} shows the bracket geometry and load case, and
\cref{fig:sbs} demonstrates placing two related figures side by side --- use
this layout for a loaded diagram next to its corresponding free-body
diagram, for example.

\begin{figure}[H]
    \centering
    \begin{tikzpicture}
        \draw[thick] (0,0) rectangle (0.3,3);
        \draw[thick, fill=tudelft-cyan!15] (0.3,1.35) rectangle (3,1.65);
        \draw[thick] (3,1.5) circle (0.12);
        \draw[->, thick] (3,1.5) -- (3,0.6) node[below] {$F(t)$};
        \draw[dashed] (0.3,1.5) -- (3,1.5);
        \draw[<->] (0.3,1.9) -- (3,1.9) node[midway, above] {$L$};
        \node[below] at (0.15,-0.1) {mounting plate};
        \node[above right] at (3,1.6) {sensor mount};
    \end{tikzpicture}
    \caption{Cantilevered sensor bracket, load case, and key dimension $L$}
    \label{fig:bracket}
\end{figure}

\begin{figure}[H]
    \centering
    \begin{subfigure}[t]{0.48\textwidth}
        \centering
        \begin{tikzpicture}
            \draw[thick,->] (0,0) -- (2,0) node[right] {$x$};
            \draw[thick,->] (0,0) -- (0,1.5) node[above] {$M(x)$};
            \draw[thick, tudelft-cyan] (0,0) -- (2,1.2);
        \end{tikzpicture}
        \caption{Bending moment diagram}
        \label{fig:sbs-moment}
    \end{subfigure}
    \hfill
    \begin{subfigure}[t]{0.48\textwidth}
        \centering
        \begin{tikzpicture}
            \draw[thick,->] (0,0) -- (2,0) node[right] {$x$};
            \draw[thick,->] (0,0) -- (0,1.5) node[above] {$V(x)$};
            \draw[thick, tudelft-cyan] (0,0) -- (2,0.9);
        \end{tikzpicture}
        \caption{Shear force diagram}
        \label{fig:sbs-shear}
    \end{subfigure}
    \caption{Internal load distributions along the bracket span}
    \label{fig:sbs}
\end{figure}

\section{Nomenclature}
\nomenclature[B]{$M$}{Bending moment \nomunit{\newton\meter}}
\nomenclature[B]{$F$}{Applied load \nomunit{\newton}}
\nomenclature[B]{$\sigma$}{Bending stress \nomunit{\mega\pascal}}
\nomenclature[G]{$\sigma_y$}{Yield stress of 6061-T6 aluminium \nomunit[276]{\mega\pascal}}
\nomenclature[A]{FoS}{Factor of safety}
\nomenclature[A]{DSE}{Design Synthesis Exercise}

This section illustrates the \texttt{questions} environment, tables,
side-by-side figures, TikZ diagrams, \texttt{siunitx} units, and the
nomenclature/list-of-symbols machinery inherited from the class. Replace it
with your own analysis chapters.
